\documentclass[11pt,a4paper]{article}

\usepackage[utf8]{inputenc}
\usepackage[T1]{fontenc}
\usepackage[spanish]{babel}
\usepackage[margin=2.5cm]{geometry}
\usepackage{graphicx}
\usepackage{booktabs}
\usepackage{array}
\usepackage{listings}
\usepackage{xcolor}
\usepackage{tikz}
\usetikzlibrary{arrows.meta, positioning, shapes.multipart}
\usepackage{hyperref}
\usepackage{enumitem}

\hypersetup{
    colorlinks=true,
    linkcolor=blue!50!black,
    urlcolor=blue!50!black
}

\lstdefinestyle{terminal}{
    basicstyle=\ttfamily\small,
    backgroundcolor=\color{black!5},
    frame=single,
    breaklines=true,
    columns=fullflexible,
    showstringspaces=false
}
\lstset{style=terminal}

\title{\textbf{Mini Red Social} \\ \large Informe Técnico \\ Algoritmos y Estructuras de Datos}
\author{
    Integrante 1 \\
    Integrante 2 \\
    Integrante 3 \\
    Integrante 4
}
\date{31 de julio de 2026}

\begin{document}

\maketitle
\tableofcontents
\newpage

% ---------------------------------------------------------------------
\section{Introducción}

Este proyecto implementa el núcleo de funcionamiento de una red social
simplificada: usuarios, relaciones de amistad, publicaciones, comentarios
y reacciones, replicando el tipo de operaciones que sostiene una red
social real (registrar personas, conectar amistades, publicar contenido y
calcular relaciones no triviales entre usuarios, como el camino de
amistad o las sugerencias de amistad).

La restricción central del enunciado es que ninguna estructura de datos
proviene de la STL de C++. \texttt{Vector}, listas enlazadas,
\texttt{Deque}, \texttt{AVLTree}, \texttt{Heap}, \texttt{HashTable} y
\texttt{Graph} están implementadas desde cero por el equipo. Sobre esas
siete estructuras se construye una capa de entidades (\texttt{User},
\texttt{Post}, \texttt{Comment}) y un servicio orquestador
(\texttt{SocialNetwork}) que expone las funcionalidades pedidas en el
enunciado a través de una interfaz de consola (CLI).

El desarrollo fue incremental: cada estructura de datos se implementó,
probó y aprobó por separado antes de pasar a la siguiente, y recién con
las siete estructuras listas se construyó la capa de dominio y el
servicio que las integra.

% ---------------------------------------------------------------------
\section{Arquitectura del sistema}

El sistema está organizado en cuatro capas, de arriba hacia abajo:

\begin{enumerate}
    \item \textbf{CLI} (\texttt{src/main.cpp}): menús jerárquicos de
    consola (usuarios, publicaciones, amistades, estadísticas /
    rendimiento, dataset).
    \item \textbf{Servicio orquestador} (\texttt{SocialNetwork}, en
    \texttt{include/servicios} y \texttt{src/servicios}): conecta todo lo
    demás y expone las operaciones del enunciado.
    \item \textbf{Entidades de dominio} (\texttt{User}, \texttt{Post},
    \texttt{Comment}, en \texttt{include/modelo} y \texttt{src/modelo}):
    representan los datos que pide el enunciado.
    \item \textbf{Estructuras de datos genéricas}
    (\texttt{include/estructuras}): \texttt{Vector}, \texttt{LinkedList},
    \texttt{Deque}, \texttt{AVLTree}, \texttt{Heap}, \texttt{HashTable} y
    \texttt{Graph}, todas implementadas desde cero.
\end{enumerate}

\texttt{SocialNetwork} mantiene internamente:

\begin{itemize}
    \item \texttt{HashTable<int, User>} y \texttt{HashTable<int, Post>}
    como índice primario por ID (búsqueda O(1) promedio).
    \item \texttt{HashTable<int, Comment>} para los comentarios,
    referenciados por ID desde cada \texttt{Post}.
    \item \texttt{AVLTree<std::string, LinkedList<int>>} como índice
    secundario por nombre. Permite nombres repetidos: cada nodo del árbol
    guarda la lista de IDs de los usuarios que comparten ese nombre.
    \item \texttt{Graph<int>} para las relaciones de amistad: camino de
    amistad (BFS), amigos en común y sugerencias.
    \item Dos \texttt{Vector<int>} (\texttt{m\_user\_ids},
    \texttt{m\_post\_ids}) que llevan el registro de IDs activos, porque
    \texttt{HashTable} no expone iteración sobre sus entradas y se
    necesita recorrer ``todos los usuarios'' o ``todas las
    publicaciones'' para calcular los rankings.
\end{itemize}

Los IDs de usuario, publicación y comentario son enteros
autoincrementales asignados por \texttt{SocialNetwork} y nunca se
reutilizan, incluso después de un borrado. Así se evita que un ID viejo
``resucite'' en el índice por nombre o en una lista de amigos que
todavía no se actualizó.

\texttt{DatasetGenerator} (en \texttt{data/generador}) puebla una
\texttt{SocialNetwork} con datos sintéticos (usuarios, amistades y
publicaciones) usando un generador congruencial lineal propio, sin
depender de \texttt{<random>}. \texttt{Benchmark} (en
\texttt{src/servicios}) mide con \texttt{<chrono>} el tiempo de las
operaciones más costosas sobre datasets de distinto tamaño y exporta los
resultados a un CSV.

% ---------------------------------------------------------------------
\section{Descripción de las estructuras utilizadas}

\begin{center}
\begin{tabular}{@{}l p{9.5cm}@{}}
\toprule
\textbf{Estructura} & \textbf{Rol en el proyecto} \\
\midrule
\texttt{Vector<T>} &
Arreglo dinámico no ordenado (capacidad duplicada al llenarse,
eliminación por swap con el último elemento). Se usa para listas de
amigos y publicaciones por usuario, resultados de búsquedas y
colecciones de IDs. \\[0.4em]
\texttt{LinkedList<T>} &
Lista simplemente enlazada. Se usa como cadena de colisiones dentro de
\texttt{HashTable} y como \emph{bucket} del índice por nombre en
\texttt{AVLTree}. \\[0.4em]
\texttt{Deque<T>} &
Cola de doble extremo con buffer circular. Se usa como cola FIFO en el
BFS de \texttt{Graph::shortest\_path}. \\[0.4em]
\texttt{AVLTree<K,V>} &
Árbol binario de búsqueda balanceado por rotaciones. Se usa como índice
secundario de usuarios por nombre. \\[0.4em]
\texttt{Heap<T>} &
Max-heap genérico con functor de comparación intercambiable. Se usa para
los rankings top-K: usuarios más activos, publicaciones con más
reacciones y sugerencias de amistad ordenadas por cantidad de amigos en
común. \\[0.4em]
\texttt{HashTable<K,V>} &
Tabla hash con manejo de colisiones por encadenamiento (usa
\texttt{LinkedList} internamente) y rehash automático al superar el
factor de carga máximo. Es el índice primario de usuarios y
publicaciones por ID, y la representación de la lista de adyacencia del
grafo. \\[0.4em]
\texttt{Graph<K>} &
Grafo no dirigido (lista de adyacencia sobre \texttt{HashTable}). Modela
la red de amistades: BFS para camino más corto, intersección de vecinos
para amigos en común, y conteo de amigos en común para sugerencias. \\
\bottomrule
\end{tabular}
\end{center}

% ---------------------------------------------------------------------
\section{Diagramas}

\subsection{Diagrama de capas}

\begin{center}
\begin{tikzpicture}[
    box/.style={draw, rounded corners, minimum width=10.5cm, minimum height=1.1cm, align=center, fill=blue!5},
    node distance=0.8cm
]
\node[box] (cli) {\textbf{CLI} (\texttt{src/main.cpp}) \\ Menús: usuarios, publicaciones, amistades, estadísticas, dataset};
\node[box, below=of cli] (service) {\textbf{Servicio orquestador}: \texttt{SocialNetwork}};
\node[box, below=of service, minimum width=5cm, xshift=-2.9cm] (model) {\textbf{Entidades} \\ \texttt{User}, \texttt{Post}, \texttt{Comment}};
\node[box, below=of service, minimum width=5cm, xshift=2.9cm] (struct) {\textbf{Estructuras de datos} \\ \footnotesize \texttt{Vector}, \texttt{LinkedList}, \texttt{Deque}, \texttt{AVLTree}, \texttt{Heap}, \texttt{HashTable}, \texttt{Graph}};

\draw[-Latex] (cli) -- (service);
\draw[-Latex] (service) -- (model);
\draw[-Latex] (service) -- (struct);
\draw[-Latex] (model) -- (struct);
\end{tikzpicture}
\end{center}

\subsection{Diagrama de clases (entidades y servicio)}

\begin{center}
\begin{tikzpicture}[
    classbox/.style={rectangle split, rectangle split parts=2, draw, align=left, minimum width=4cm, font=\footnotesize},
    node distance=0.9cm and 0.7cm
]

\node[classbox] (user) {
    \textbf{User}
    \nodepart{second}
    id, name, email \\
    registration\_date \\
    friend\_ids: Vector<int> \\
    post\_ids: Vector<int> \\
    follower\_count \\
    reactions\_received
};

\node[classbox, right=of user] (post) {
    \textbf{Post}
    \nodepart{second}
    id, owner\_id, date, text \\
    like\_count \\
    comment\_ids: Vector<int>
};

\node[classbox, right=of post] (comment) {
    \textbf{Comment}
    \nodepart{second}
    id, author\_id, post\_id \\
    text, date
};

\node[classbox, below=1.4cm of post, minimum width=8.5cm] (network) {
    \textbf{SocialNetwork}
    \nodepart{second}
    m\_users: HashTable<int,User> \\
    m\_posts: HashTable<int,Post> \\
    m\_comments: HashTable<int,Comment> \\
    m\_users\_by\_name: AVLTree<string,LinkedList<int>> \\
    m\_friendships: Graph<int>
};

\draw[-Latex] (network) -- (user);
\draw[-Latex] (network) -- (post);
\draw[-Latex] (network) -- (comment);
\end{tikzpicture}
\end{center}

\subsection{Grafo de amistades y camino más corto (BFS)}

Ejemplo real tomado de \texttt{tests/graph\_test.cpp}: un grafo con
aristas $(1,2)$, $(1,3)$, $(2,4)$, $(3,4)$ y $(4,5)$. El camino de
amistad calculado por \texttt{shortest\_path(1, 5)} (resaltado en rojo)
es $1 \rightarrow 2 \rightarrow 4 \rightarrow 5$.

\begin{center}
\begin{tikzpicture}[node/.style={draw, circle, minimum size=0.9cm}, >=Latex]
\node[node] (n1) at (0,0)   {1};
\node[node] (n2) at (2,1)   {2};
\node[node] (n3) at (2,-1)  {3};
\node[node] (n4) at (4,0)   {4};
\node[node] (n5) at (6,0)   {5};

\draw (n1) -- (n3);
\draw (n3) -- (n4);

\draw[very thick, red] (n1) -- (n2);
\draw[very thick, red] (n2) -- (n4);
\draw[very thick, red] (n4) -- (n5);
\end{tikzpicture}
\end{center}

Este mismo grafo se usa en el informe para ilustrar \texttt{common\_friends(2, 3)}:
ambos comparten como vecinos a $1$ y a $4$, por lo que el resultado es
$\{1, 4\}$ (verificado en la sección de capturas).

% ---------------------------------------------------------------------
\section{Complejidad computacional}

Sea $n$ el número de usuarios y $d$ el grado promedio (cantidad de
amigos) de un usuario.

\begin{center}
\begin{tabular}{@{}l l p{6.3cm}@{}}
\toprule
\textbf{Operación} & \textbf{Complejidad} & \textbf{Notas} \\
\midrule
\texttt{Vector::add\_item} & O(1) amortizado & duplica capacidad al llenarse \\
\texttt{Vector::remove\_index} & O(1) & swap con el último elemento \\
\texttt{Vector::search} & O(n) & lineal \\
\texttt{LinkedList::add\_last} / \texttt{add\_first} & O(1) & mantiene puntero a la cola \\
\texttt{LinkedList::search} / \texttt{remove\_item} & O(n) & lineal \\
\texttt{Deque::add\_front} / \texttt{add\_back} & O(1) amortizado & buffer circular con resize \\
\texttt{AVLTree::insert} / \texttt{search} / \texttt{remove} & O(log n) & balanceado por rotaciones \\
\texttt{Heap::insert} / \texttt{extract\_top} & O(log k) & $k$ = elementos en el heap \\
\texttt{HashTable::put} / \texttt{search} / \texttt{remove} & O(1) promedio, O(n) peor caso & depende de la calidad del hash y el factor de carga \\
\texttt{SocialNetwork::register\_user} & O(1) promedio & inserción en \texttt{HashTable} + \texttt{AVLTree} (O(log u), $u$ = nombres distintos) \\
\texttt{SocialNetwork::add\_friend} / \texttt{remove\_friend} & O(d) & recorre la lista de adyacencia para evitar aristas duplicadas \\
\texttt{SocialNetwork::friendship\_path} (BFS) & O(n + a) & $a$ = número de aristas del grafo \\
\texttt{SocialNetwork::common\_friends} & O(d) & intersección lineal de dos listas de adyacencia \\
\texttt{SocialNetwork::friend\_suggestions} & O(n $\cdot$ d) & recorre todos los vértices y, para cada uno, su lista de adyacencia \\
\texttt{SocialNetwork::most\_active\_users} / \texttt{most\_reacted\_posts} & O(n log n) & inserta todos los elementos en un \texttt{Heap} y extrae el top-k \\
\bottomrule
\end{tabular}
\end{center}

El punto más débil identificado es \texttt{friend\_suggestions}: la
implementación actual de \texttt{Graph} recorre \textbf{todos} los
vértices del grafo por cada consulta (no solo los amigos de los amigos),
lo que la vuelve O(n $\cdot$ d) en vez de O(d\textsuperscript{2}). Esto se
confirma en los resultados experimentales de la siguiente sección y
queda documentado aquí como una limitación conocida, no como un error
oculto.

% ---------------------------------------------------------------------
\section{Resultados experimentales}

Los tiempos se midieron con \texttt{Benchmark::run}, que genera un
dataset sintético con \texttt{DatasetGenerator} (8 amigos promedio y 3
publicaciones promedio por usuario) y mide con \texttt{<chrono>} seis
operaciones sobre tres tamaños de dataset. Ejecutado en la máquina de
desarrollo del equipo:

\begin{center}
\small
\begin{tabular}{@{}r r r r r r r@{}}
\toprule
\textbf{Usuarios} & \textbf{Generación} & \textbf{Camino de} & \textbf{Amigos en} & \textbf{Sugerencias} & \textbf{Usuarios más} & \textbf{Posts más} \\
 & \textbf{(ms)} & \textbf{amistad (ms)} & \textbf{común (ms)} & \textbf{(ms)} & \textbf{activos (ms)} & \textbf{reaccionados (ms)} \\
\midrule
1\,000  & 27.76   & 0.38 & 0.11 & 26.13  & 0.13 & 0.40 \\
5\,000  & 500.07  & 2.06 & 0.24 & 304.95 & 0.17 & 0.71 \\
10\,000 & 1\,898.13 & 3.59 & 0.31 & 703.10 & 0.34 & 1.25 \\
\bottomrule
\end{tabular}
\end{center}

Observaciones:

\begin{itemize}
    \item El \textbf{camino de amistad (BFS)} y los \textbf{rankings
    top-K} escalan bien: crecen mucho más lento que el tamaño del
    dataset, consistente con sus complejidades O(n + a) y O(n log n).
    \item Las \textbf{sugerencias de amistad} son, con claridad, la
    operación más costosa y la que peor escala (crece casi proporcional
    al número de usuarios), confirmando el análisis de la sección
    anterior.
    \item La \textbf{generación del dataset} domina el tiempo total
    porque cada \texttt{add\_friend}/\texttt{create\_post} pasa por
    varias estructuras (\texttt{HashTable}, \texttt{AVLTree},
    \texttt{Graph}); es esperable y no es una operación que el usuario
    final repita en un flujo normal de uso.
\end{itemize}

Estos números se regeneran en
\texttt{data/generador/benchmark\_results.csv} al ejecutar la opción
``Ejecutar benchmark'' del menú de Estadísticas del CLI.

% ---------------------------------------------------------------------
\section{Capturas del funcionamiento}

Salida real del CLI (\texttt{mini\_socialnetwork}) generando un dataset
pequeño y consultando el ranking de usuarios más activos:

\begin{lstlisting}
--- Dataset ---
1. Generar dataset sintetico
0. Volver
Opcion: 1
Cantidad de usuarios a generar: 50
Amigos promedio por usuario: 5
Publicaciones promedio por usuario: 2
Dataset generado. Usuarios totales: 50, publicaciones totales: 100

--- Estadisticas / Rendimiento ---
1. Usuarios mas activos
2. Publicaciones con mas reacciones
3. Ejecutar benchmark (exporta CSV)
0. Volver
Opcion: 1
Cuantos mostrar: 5
ID    Nombre                  Correo                      Amigos   Posts
1     Pedro Flores            user1@mail.com              15       2
50    Ana Rios                user50@mail.com             5        2
49    Jose Vega               user49@mail.com             5        2
48    Sofia Gomez             user48@mail.com             5        2
47    Diego Torres            user47@mail.com             5        2
\end{lstlisting}

Perfil de usuario formateado:

\begin{lstlisting}
----------------------------------------
 Perfil de Pedro Flores
----------------------------------------
 ID:              1
 Correo:          user1@mail.com
 Registrado:      2026-06-25
 Amigos:          15
 Seguidores:      0
 Publicaciones:   2
 Reacciones recibidas: 3
----------------------------------------
\end{lstlisting}

Salida de \texttt{graph\_test}, sobre el grafo de la sección de
diagramas (confirma el camino BFS $1 \to 2 \to 4 \to 5$ y los amigos en
común $\{1, 4\}$ entre los usuarios 2 y 3):

\begin{lstlisting}
=== Prueba de Graph<int> ===
Vertices: 5
Tiene arista 1-2: si
Camino 1->5: 1 2 4 5
Amigos en comun 2 y 3: 1 4
Sugerencias para 1: 4
Tiene arista 1-2 tras borrar: no
\end{lstlisting}

% ---------------------------------------------------------------------
\section{Conclusiones}

\begin{itemize}
    \item Fue posible construir toda la capa de datos de una red social
    (usuarios, amistades, publicaciones, comentarios, reacciones) usando
    únicamente estructuras propias, sin recurrir a la STL.
    \item Separar el proyecto en estructuras genéricas
    (\texttt{include/estructuras}), entidades de dominio
    (\texttt{include/modelo}) y un servicio orquestador
    (\texttt{include/servicios}) permitió que cada estructura se probara
    de forma aislada antes de integrarla, y que el servicio no dependiera
    de los detalles internos de cada estructura.
    \item El cuello de botella real del sistema no está en las
    estructuras base (todas se comportan según su complejidad teórica)
    sino en un algoritmo de más alto nivel
    (\texttt{friend\_suggestions}), lo cual es el tipo de hallazgo que
    solo aparece al medir con datasets de tamaño creciente en vez de
    asumir la complejidad sobre el papel.
    \item Queda como trabajo futuro, fuera del alcance de este proyecto:
    optimizar \texttt{friend\_suggestions} para que recorra solo
    amigos-de-amigos en vez de todos los vértices, y probar el sistema
    con datasets de cientos de miles de usuarios (el diseño está pensado
    para escalar a ese tamaño, pero se validó experimentalmente hasta
    10\,000 usuarios por el tiempo disponible para el proyecto).
\end{itemize}

\end{document}
